\chapter{Requirementsanalyse}%
\label{ch:requirementsanalyse}

In de requirementsanalyse worden twee vragen beantwoord. Ten eerste wordt er bepaald welke use cases of functies zullen worden uitgevoerd in de PoC's om Python en Rexx aan elkaar te vergelijken. Ten tweede wordt er bepaald aan de hand van welke succescriteria zal worden bepaald of een use case succesvol kan worden uitgevoerd door een van beide talen.

Er worden verschillende methodes toegepast om deze requirements te verkrijgen. Er worden interviews uitgevoerd, en daarnaast wordt een enquête opgesteld, om zo input te krijgen van de mainframe gemeenschap.

\section{Interviews}

\section{Enquête}
De bedoeling van de enquête was om zoveel mogelijk input te verkrijgen van de mainframe gemeenschap, omdat de resultaten van dit onderzoek mede bedoeld zijn voor die doelgroep. Daarnaast is de gemeenschap zodanig klein en behulpzaam dat dit een goede manier leek om discussies gaande te krijgen en zo nog meer input te verkrijgen. De enquête is op het moment van schrijven en verwerking van de resultaten 21 keer ingevuld. Daarnaast is er ook veel reactie gekomen op de Linkedin post om de enquête te promoten. Al die reacties samen hebben ervoor gezorgd dat er meer dan voldoende requirements en use cases zijn verzameld om uit te kiezen voor dit onderzoek.

\section{Selectie}

Zoals vermeld zijn er vrij veel requirements en use cases verzameld om dit onderzoek uit te voeren. Uit deze verzameling is een selectie gemaakt om verder te gaan. Deze selectie is gebeurd op basis van verschillende criteria. Voor zowel de use cases als requirements geldt dat het mogelijk moet zijn om uit te voeren zonder toegang tot een productie systeem. Dit zorgt ervoor dat alles te maken met CPU consumptie, snelheid of algemene performantie niet mogelijk is. Daarnaast zijn subjectieve succescriteria, en use cases die buiten de scope van de huidige vergelijking vallen, ook niet meegenomen in de laatste selectie. Hieronder wordt een overzicht gegeven van de use cases en requirements die zullen worden gebruikt in het hoofdstuk van proof-of-concept.
